\documentclass[twocolumn,10pt]{article}

% ============================================================================
% Packages
% ============================================================================
\usepackage[utf8]{inputenc}
\usepackage[T1]{fontenc}
\usepackage{amsmath,amssymb,amsthm}
\usepackage{mathtools}
\usepackage{physics}
\usepackage{graphicx}
\usepackage[margin=1.8cm,top=2.2cm,bottom=2.2cm]{geometry}
\usepackage{natbib}
\usepackage{hyperref}
\usepackage{cleveref}
\usepackage{booktabs}
\usepackage{multirow}
\usepackage{array}
\usepackage{float}
\usepackage{caption}
\usepackage{enumitem}
\usepackage{xcolor}
\usepackage[expansion=false]{microtype}
\usepackage{bm}
\usepackage{tabularx}
\usepackage{algorithm}
\usepackage{algpseudocode}

% ============================================================================
% Theorem environments
% ============================================================================
\newtheorem{theorem}{Theorem}[section]
\newtheorem{lemma}[theorem]{Lemma}
\newtheorem{proposition}[theorem]{Proposition}
\newtheorem{corollary}[theorem]{Corollary}
\newtheorem{definition}[theorem]{Definition}
\newtheorem{axiom}{Axiom}
\newtheorem{principle}[theorem]{Principle}

\theoremstyle{remark}
\newtheorem{remark}[theorem]{Remark}

% ============================================================================
% Custom commands
% ============================================================================
\newcommand{\Sk}{S_{\mathrm{k}}}
\newcommand{\St}{S_{\mathrm{t}}}
\newcommand{\Se}{S_{\mathrm{e}}}
\newcommand{\Sosc}{S_{\mathrm{osc}}}
\newcommand{\Scat}{S_{\mathrm{cat}}}
\newcommand{\Spart}{S_{\mathrm{part}}}
\newcommand{\kB}{k_{\mathrm{B}}}
\newcommand{\dcat}{d_{\mathrm{cat}}}
\newcommand{\Hplus}{\mathrm{H}^{+}}
\newcommand{\Otwo}{\mathrm{O}_{2}}
\newcommand{\Rorder}{R}
\newcommand{\PLV}{\mathrm{PLV}}
\newcommand{\Mstate}{\mathcal{M}}
\newcommand{\Ccat}{\mathcal{C}}
\newcommand{\Sspace}{\mathcal{S}}
\newcommand{\Gstate}{\Gamma}
\newcommand{\Pdecay}{P_{\mathrm{decay}}}
\newcommand{\Tdecay}{T_{\mathrm{decay}}}
\newcommand{\Hfield}{\mathcal{H}^+}
\newcommand{\partition}[4]{(#1,\,#2,\,#3,\,#4)}
\DeclareMathOperator{\sgn}{sgn}

% ============================================================================
% Title
% ============================================================================
\title{%
  \textbf{Pharmacological Equations of State for Hybrid Microfluidic Circuits: Drug Action as Regime Transition in Bounded Neural Phase Space}
}

\author{
  Kundai~F.~Sachikonye%
  \thanks{Corresponding author. Email: kundai.sachikonye@tum.de}
  \\[4pt]
  \small Technical University of Munich \\
  \small Munich, Germany
}

\date{}

% ============================================================================
\begin{document}
% ============================================================================
\twocolumn[
  \maketitle
  \begin{@twocolumnfalse}
  \begin{abstract}
\noindent
We derive pharmacological equations of state for hybrid microfluidic circuits---bounded neural phase space systems integrating oscillatory, categorical, and partition dynamics---from three axioms: bounded phase space, the No Null State principle, and finite observational resolution. Drug action is formulated as structural factor modification in the universal equation of state $PV = N\kB T \cdot \mathcal{S}(V, N, \{n_i, \ell_i, m_i, s_i\})$, where the temperature-independent structural factor $\mathcal{S}$ encodes partition geometry and is the sole target of pharmacological intervention. Central to the framework is the \emph{categorical aperture}: a topological constraint on phase space that achieves molecular selectivity at rigorously zero thermodynamic work ($W_{\mathrm{aperture}} = 0$), operating through electromagnetic field configuration rather than information acquisition, thereby resolving the Maxwell demon paradox without invoking Landauer erasure costs.

We classify drug action through three mechanisms: (i)~aperture installation---the drug molecule establishes a new topological constraint on neural phase space, classified by electromagnetic multipole order (monopole, dipole, quadrupole); (ii)~coupling strength modulation---the drug alters the Kuramoto coupling parameter $K$, shifting the system toward or away from the critical synchronization threshold $K_c = 2\sigma(\omega)/\langle k \rangle$; (iii)~regime transition induction---the drug drives the circuit between five operational regimes (coherent, turbulent, hierarchical cascade, aperture-dominated, phase-locked), each possessing a distinct equation of state. Drug classes correspond to characteristic regime transition pathways: SSRIs traverse Turbulent $\to$ Aperture-Dominated $\to$ Phase-Locked $\to$ Coherent; SNRIs traverse Turbulent $\to$ Hierarchical Cascade $\to$ Coherent; tricyclic antidepressants traverse Turbulent $\to$ Phase-Locked $\to$ Coherent.

The framework predicts cross-modal therapeutic equivalence: chemically diverse drugs that induce the same regime transition produce equivalent clinical outcomes despite binding profile distances of $5.8$--$6.7$ in partition coordinate space, explaining the convergent $\sim\!60\%$ response rate across antidepressant classes. Variance minimization $F = \kB T \cdot \sigma^2(\phi)$ provides the thermodynamic basis for therapeutic action---drugs reduce phase variance toward the thermal floor $\sigma^2_{\min} = \kB T / K_{\mathrm{coupling}}$. The 2--4 week antidepressant onset delay is derived as the time required for sustained coupling enhancement to cross $K_c$.

We extend the Neural Partition Language (NPL) with pharmacological primitives---aperture operations, regime transitions, frequency matching, and trajectory modification operators---grounded in the equations of state. The resulting pharmacological NPL (pNPL) expresses drug action as declarative constraint satisfaction: specify the target regime, and the partition structure determines the required trajectory. Experimental validation pathways include longitudinal EEG phase-locking value monitoring during antidepressant treatment, dose--response curves for structural factor modification, and real-time regime classification for pharmacological titration.

\vspace{6pt}
\noindent\textbf{Keywords:} pharmacological equations of state, categorical aperture, regime transition, Kuramoto synchronization, structural factor, variance minimization, neural partition language, drug action, phase-locking value, bounded phase space
  \end{abstract}
  \vspace{12pt}
  \end{@twocolumnfalse}
]

% ============================================================================
\section{Introduction}
\label{sec:intro}
% ============================================================================

\subsection{The Pharmacological Problem}

Contemporary pharmacology rests on a binding paradigm: drug efficacy is predicted from receptor affinity ($K_i$), selectivity ratios, and dose--response curves~\citep{kenakin2004principles}. This framework has produced effective therapeutics but fails to explain three persistent anomalies. First, chemically unrelated drugs---selective serotonin reuptake inhibitors (SSRIs), tricyclic antidepressants (TCAs), and atypical agents---produce statistically indistinguishable clinical response rates of approximately $60\%$ for major depression~\citep{cipriani2018comparative}. Second, drugs with high receptor affinity do not uniformly produce strong clinical effects, while promiscuously binding agents sometimes outperform selective ones~\citep{roth2004multiplicity}. Third, antidepressant receptor occupancy occurs within hours, yet clinical response requires 2--4 weeks~\citep{stahl2021essential}.

These anomalies suggest that the relevant pharmacological variable is not receptor binding per se but the \emph{dynamical state} of the neural circuit being perturbed. We demonstrate that drug action is rigorously described as structural factor modification in the thermodynamic equation of state of the neural circuit, with therapeutic response corresponding to regime transition in bounded phase space.

\subsection{The Hybrid Microfluidic Circuit}

The neural system is a hybrid microfluidic circuit: a bounded physical system integrating fluid transport (cerebrospinal and interstitial flow, ionic currents), electromagnetic coupling (synaptic transmission, field potentials), and geometric partitioning (dendritic arborization, cortical columns, nuclear boundaries)~\citep{whitesides2006origins,stone2004engineering}. These three modalities---oscillatory dynamics, categorical completion, and partition geometry---operate simultaneously and are governed by a single thermodynamic framework.

The circuit occupies bounded phase space (finite energy, finite spatial extent, finite number of neurons), admits hierarchical partition into distinguishable categories (cortical columns, frequency bands, functional networks), and satisfies the No Null State principle (neural activity never ceases; even ``resting state'' is a specific dynamical configuration). From these three constraints alone, we derive complete equations of state governing circuit dynamics and, consequently, governing how drugs perturb those dynamics.

\subsection{Scope and Organization}

This paper is organized as follows. Section~\ref{sec:axioms} establishes the three foundational axioms and derives the triple equivalence governing hybrid circuit thermodynamics. Section~\ref{sec:apertures} defines categorical apertures, proves zero-work operation, and classifies drug molecules by aperture type. Section~\ref{sec:regimes} derives the five operational regimes with their equations of state. Section~\ref{sec:drug_action} formulates drug action as structural factor modification and derives regime transition pathways for major drug classes. Section~\ref{sec:frequency} establishes the frequency matching and timescale hierarchy governing drug--circuit coupling. Section~\ref{sec:trajectories} develops pharmacological trajectories in S-entropy space and the variance minimization principle. Section~\ref{sec:pnpl} extends the Neural Partition Language with pharmacological primitives. Section~\ref{sec:time} derives temporal dynamics of drug action from circuit completion mechanics. Section~\ref{sec:predictions} presents experimental predictions and validation pathways. Section~\ref{sec:clinical} derives clinical implications. Section~\ref{sec:discussion} discusses the framework in relation to existing pharmacological theory.

% ============================================================================
\section{Foundational Axioms and the Triple Equivalence}
\label{sec:axioms}
% ============================================================================

We begin from three axioms that constrain the class of physical systems under consideration. These axioms are not specific to neuroscience---they apply to any bounded physical system---but their consequences for neural circuits are immediate and far-reaching.

\begin{axiom}[Bounded Phase Space]
\label{ax:bounded}
Every physical system occupies a phase space $\Omega$ of finite measure: $\mu(\Omega) < \infty$.
\end{axiom}

For a neural circuit, this reflects physical reality: the brain has finite volume ($\sim\!1400$ cm$^3$), finite energy consumption ($\sim\!20$ W), and a finite number of neurons ($\sim\!10^{11}$). The accessible phase space, while vast, is bounded.

\begin{axiom}[No Null State]
\label{ax:nonull}
A system must occupy exactly one categorical state at every instant. There is no ``null'' state, no ``between'' state, and no moment at which the system fails to belong to a category.
\end{axiom}

This axiom has a profound consequence: \emph{oscillation arises from categorical necessity, not from dynamical forces}~\citep{buzsaki2006rhythms}. In a bounded phase space with finitely many categories, a system that cannot occupy a null state must perpetually transition between categories. Forces provide the \emph{mechanism} for transitions; categorical necessity provides the \emph{reason}.

\begin{axiom}[Finite Observational Resolution]
\label{ax:resolution}
Any observation of the system has finite resolution $\delta > 0$. States separated by less than $\delta$ are observationally indistinguishable.
\end{axiom}

Axiom~\ref{ax:resolution} ensures that the number of distinguishable states is finite: $N = \mu(\Omega)/\delta^d$, where $d$ is the phase space dimensionality.

\subsection{Partition Coordinates}

From Axioms~\ref{ax:bounded} and \ref{ax:resolution}, categorical partitioning of bounded spherical phase space generates four discrete coordinates.

\begin{definition}[Partition Coordinates]
\label{def:partition_coords}
The state of a system at partition depth $n$ is specified by
\begin{equation}
  \mathbf{q} = \partition{n}{\ell}{m}{s},
\end{equation}
where $n \in \mathbb{N}^+$ (depth level), $\ell \in \{0, 1, \ldots, n-1\}$ (complexity), $m \in \{-\ell, \ldots, +\ell\}$ (orientation), and $s \in \{-\frac{1}{2}, +\frac{1}{2}\}$ (chirality).
\end{definition}

\begin{theorem}[Shell Capacity]
\label{thm:capacity}
The capacity at depth $n$ is $C(n) = 2n^2$.
\end{theorem}

\begin{proof}
For depth $n$, complexity ranges $\ell \in \{0, \ldots, n-1\}$. Each $\ell$ level contains $2\ell + 1$ orientation states, each with 2 chirality states:
\begin{equation}
  C(n) = \sum_{\ell=0}^{n-1} 2(2\ell + 1) = 2n^2.
\end{equation}
\end{proof}

\subsection{The Triple Equivalence}

\begin{definition}[Partition Depth]
\label{def:depth}
The partition depth of a system with $K$ hierarchical levels, each containing $k_i$ categories, is
\begin{equation}
  M = \sum_{i=1}^{K} \log_b k_i,
\end{equation}
where $b$ is the partition basis ($b = 3$ is optimal among integers).
\end{definition}

\begin{theorem}[Triple Equivalence]
\label{thm:triple}
For a system satisfying Axioms~\ref{ax:bounded}--\ref{ax:resolution}, three entropy formulations are identical:
\begin{equation}
  \boxed{\Sosc = \Scat = \Spart = \kB M \ln n}
  \label{eq:triple}
\end{equation}
where $\Sosc$ is oscillatory entropy (phase configuration), $\Scat$ is categorical entropy (state distribution), and $\Spart$ is partition entropy (hierarchical structure).
\end{theorem}

\begin{proof}
By Axiom~\ref{ax:bounded}, the phase space is bounded with finite microstate count $N$. By Axiom~\ref{ax:resolution}, distinguishable states number $N = \mu(\Omega)/\delta^d$. At partition depth $M$ with branching ratio $n$, the system has $n^M$ categories.

For the partition description: $\Spart = \kB \ln(n^M) = \kB M \ln n$.

For the categorical description: by the No Null State axiom, exactly one category is occupied at each instant. Maximum entropy over $n^M$ equiprobable categories gives $\Scat = \kB \ln(n^M) = \kB M \ln n$.

For the oscillatory description: each categorical state corresponds to a unique phase configuration of the coupled oscillator network~\citep{kuramoto1984chemical}. The number of distinct configurations at resolution $\delta$ equals $n^M$, giving $\Sosc = \kB \ln(n^M) = \kB M \ln n$.

Since all three descriptions enumerate the same microstates, the entropies are identical.
\end{proof}

\subsection{S-Entropy Coordinate Space}

\begin{definition}[S-Entropy Coordinates]
\label{def:sentropy}
The S-entropy coordinate system maps the partition state to the unit cube:
\begin{equation}
  \mathbf{S} = (\Sk, \St, \Se) \in [0, 1]^3,
\end{equation}
where $\Sk$ is knowledge entropy (configuration uncertainty), $\St$ is temporal entropy (timing uncertainty), and $\Se$ is evolution entropy (trajectory uncertainty).
\end{definition}

Navigation in S-entropy space proceeds by ternary trisection with complexity $\mathcal{O}(\log_3 N)$, yielding a $37\%$ reduction in navigational steps compared to binary search.

\begin{theorem}[Temperature Factorization]
\label{thm:temp}
For any observable $\mathcal{O}$ of a partition-based system,
\begin{equation}
  \mathcal{O} = (\kB T) \cdot F(M, n),
\end{equation}
where $F(M, n)$ depends on partition geometry but not on temperature.
\end{theorem}

This factorization is critical for pharmacology: it means that drugs modify the temperature-independent structural factor $F$---the geometry of the partition---not the thermal energy scale. Drug action is geometric, not thermal.

% ============================================================================
\section{Categorical Apertures as Drug Mechanism}
\label{sec:apertures}
% ============================================================================

\subsection{Definition and Zero-Work Theorem}

\begin{definition}[Categorical Aperture]
\label{def:aperture}
A categorical aperture $\mathcal{A}$ is a topological constraint on phase space that restricts accessible states based on position in configuration space, not on dynamical variables:
\begin{equation}
  \mathcal{A}: \Omega \to \Omega' \subseteq \Omega, \quad \Omega' = \{(q, p) \in \Omega : q \in \mathcal{D}\},
\end{equation}
where $\mathcal{D} \subseteq \mathcal{Q}$ is a domain in configuration space.
\end{definition}

\begin{theorem}[Zero-Work Aperture Filtering]
\label{thm:zero_work}
The thermodynamic work performed by a categorical aperture is identically zero:
\begin{equation}
  \boxed{W_{\mathrm{aperture}} = 0}
\end{equation}
\end{theorem}

\begin{proof}
Work is $W = \int \mathbf{F} \cdot d\mathbf{q}$. The aperture constrains position ($q \in \mathcal{D}$) but applies no force. The Hamiltonian with aperture is $H_{\mathrm{aperture}}(q,p) = H_0(q,p) + V_{\mathrm{aperture}}(q)$, where
\begin{equation}
  V_{\mathrm{aperture}}(q) = \begin{cases} 0 & q \in \mathcal{D}, \\ +\infty & q \notin \mathcal{D}. \end{cases}
\end{equation}
For particles within $\mathcal{D}$, $V = 0$ and $\nabla_q V = 0$: no work performed. For particles outside $\mathcal{D}$, the infinite potential prevents entry without momentum transfer. Thus $W_{\mathrm{aperture}} = 0$ identically.
\end{proof}

\begin{corollary}[Resolution of the Maxwell Demon Paradox]
\label{cor:demon}
Since $W_{\mathrm{aperture}} = 0$, categorical aperture filtering does not constitute a measurement in the Szil\'ard--Landauer sense~\citep{szilard1929entropieverminderung,landauer1961irreversibility,bennett1982thermodynamics}. No information is acquired, no memory is filled, and Landauer's erasure bound $W \geq \kB T \ln 2$ does not apply. The biological selectivity achieved by ion channels~\citep{doyle1998structure,hille2001ion}, enzyme active sites, and drug--receptor interactions is accomplished through topological constraint, not through information processing.
\end{corollary}

\subsection{Drug Molecules as Exogenous Apertures}

Drug molecules, upon entering the neural circuit, establish new topological constraints on the circuit's phase space. The drug does not ``measure'' the circuit state; it imposes a geometric boundary condition.

\begin{definition}[Pharmacological Aperture]
\label{def:pharm_aperture}
A pharmacological aperture is a categorical aperture installed by an exogenous molecule (drug) that restricts the neural circuit's accessible phase space:
\begin{equation}
  \mathcal{A}_{\mathrm{drug}}: \Omega_{\mathrm{neural}} \to \Omega'_{\mathrm{neural}} \subset \Omega_{\mathrm{neural}}.
\end{equation}
The restriction ratio $\alpha = \mu(\Omega_{\mathrm{neural}})/\mu(\Omega'_{\mathrm{neural}})$ determines the aperture's catalytic power.
\end{definition}

\subsection{Aperture Taxonomy by Multipole Order}

Categorical apertures in biological systems are electromagnetic field configurations that select molecules by charge interaction. We classify pharmacological apertures by the multipole order of the selecting field.

\begin{definition}[Monopole Aperture]
\label{def:monopole}
A monopole aperture generates a radial electric field $\mathbf{E} \propto \hat{r}/r^2$ selecting by net charge. Pharmacological realization: drugs binding a single transporter target (e.g., SSRIs at SERT with $K_i < 10$ nM at one site).
\end{definition}

\begin{definition}[Dipole Aperture]
\label{def:dipole}
A dipole aperture generates a two-lobed field $\mathbf{E} \propto (2\cos\theta\,\hat{r} + \sin\theta\,\hat{\theta})/r^3$ selecting by charge distribution asymmetry. Pharmacological realization: drugs binding two transporter targets (e.g., SNRIs at SERT and NET).
\end{definition}

\begin{definition}[Quadrupole Aperture]
\label{def:quadrupole}
A quadrupole aperture generates a four-lobed field $\mathbf{E} \propto (3\cos^2\theta - 1)/r^4$ selecting by quadrupole moment. Pharmacological realization: drugs with three or more high-affinity targets (e.g., TCAs binding SERT, NET, muscarinic, histaminergic, and adrenergic receptors).
\end{definition}

\begin{theorem}[Aperture Catalytic Factor]
\label{thm:catalytic}
An aperture restricting accessible phase space by factor $\alpha$ yields effective catalytic enhancement:
\begin{equation}
  k_{\mathrm{eff}} = \alpha \cdot k_{\mathrm{uncatalyzed}},
\end{equation}
with typical pharmacological values $\alpha \sim 10^5$--$10^{11}$.
\end{theorem}

The catalytic factor scales with the number of high-affinity targets: monopole apertures achieve $\alpha \sim 10^5$, dipole $\sim 10^7$, quadrupole $\sim 10^{11}$. This explains why broadly binding drugs (TCAs) often show stronger acute neurophysiological effects than selective agents (SSRIs), at the cost of increased off-target effects.

% ============================================================================
\section{Five Operational Regimes and Their Equations of State}
\label{sec:regimes}
% ============================================================================

The dynamics of a neural circuit satisfying Axioms~\ref{ax:bounded}--\ref{ax:resolution} fall into five distinct operational regimes, each characterized by specific values of the Kuramoto order parameter $\Rorder$ and the phase variance $\sigma^2(\phi)$~\citep{kuramoto1984chemical,strogatz2000kuramoto}.

\begin{definition}[Kuramoto Order Parameter]
\label{def:kuramoto}
For $N$ coupled oscillators with phases $\{\phi_i\}_{i=1}^N$, the order parameter is
\begin{equation}
  \Rorder \, e^{i\psi} = \frac{1}{N} \sum_{j=1}^{N} e^{i\phi_j},
\end{equation}
where $\Rorder \in [0, 1]$ measures phase coherence and $\psi$ is the mean phase.
\end{definition}

\begin{theorem}[Universal Equation of State]
\label{thm:eos}
Each operational regime possesses an equation of state of the universal form
\begin{equation}
  \boxed{PV = N\kB T \cdot \mathcal{S}(V, N, \{n_i, \ell_i, m_i, s_i\})}
  \label{eq:eos}
\end{equation}
where the structural factor $\mathcal{S}$ depends on partition geometry but is independent of temperature.
\end{theorem}

The five regimes and their structural factors are:

\paragraph{Regime I: Coherent ($\Rorder > 0.8$).}
All oscillators are phase-locked, hierarchical scales are fully active, and the system exhibits maximum information throughput. This is the regime of healthy waking cognition.
\begin{equation}
  \mathcal{S}_{\mathrm{coherent}} = 1 + \frac{\Rorder^2}{1 - \Rorder^2}.
\end{equation}
For $\Rorder = 0.95$ (healthy resting state), $\mathcal{S}_{\mathrm{coherent}} = 10.26$.

\paragraph{Regime II: Turbulent ($\Rorder < 0.3$).}
Phase coherence is lost, phase variance is high, hierarchical depth collapses. This is the regime of major depression~\citep{fingelkurts2006impaired,leuchter2012resting}.
\begin{equation}
  \mathcal{S}_{\mathrm{turbulent}} = 1 - \frac{\sigma^2(\phi)}{2\pi^2}.
\end{equation}
For $\Rorder = 0.31$ (depressed state), $\sigma^2 \approx 0.69$, giving $\mathcal{S}_{\mathrm{turbulent}} = 0.965$.

\paragraph{Regime III: Hierarchical Cascade.}
Multi-scale information flow with intermediate coherence; some scales exhibit cascade failure.
\begin{equation}
  \mathcal{S}_{\mathrm{cascade}} = \prod_{i} \left(1 + \frac{F_i^{\mathrm{out}}}{F_i^{\mathrm{in}}}\right),
\end{equation}
where $F_i^{\mathrm{in/out}}$ are information fluxes at hierarchical level $i$.

\paragraph{Regime IV: Aperture-Dominated.}
Geometric aperture filtering dominates dynamics; effective pressure scales quadratically with partition depth.
\begin{equation}
  \mathcal{S}_{\mathrm{aperture}} = \frac{n^2}{n_{\max}^2}.
\end{equation}

\paragraph{Regime V: Phase-Locked Network.}
Synchronization dominates; coupling strength exceeds the natural frequency spread.
\begin{equation}
  \mathcal{S}_{\mathrm{sync}} = 1 + \frac{K_{\mathrm{coupling}}}{\sigma(\omega)}.
\end{equation}

\begin{theorem}[Critical Coupling Threshold]
\label{thm:critical}
The transition from turbulent to phase-locked regime occurs at
\begin{equation}
  \boxed{K_c = \frac{2\sigma(\omega)}{\langle k \rangle}}
  \label{eq:kc}
\end{equation}
where $\langle k \rangle$ is the mean degree of the coupling network and $\sigma(\omega)$ is the standard deviation of the natural frequency distribution~\citep{strogatz2000kuramoto,acebron2005kuramoto}.
\end{theorem}

The critical coupling $K_c$ is the central quantity of pharmacology in this framework: it is the threshold a drug must help the circuit cross to achieve therapeutic response.

% ============================================================================
\section{Drug Action as Structural Factor Modification}
\label{sec:drug_action}
% ============================================================================

\subsection{The Pharmacological Equation of State}

\begin{theorem}[Pharmacological Structural Factor]
\label{thm:drug_structural}
Drug action modifies the structural factor while preserving the universal equation of state:
\begin{equation}
  PV = N\kB T \cdot \mathcal{S}_{\mathrm{drug}}(V, N, \{n_i, \ell_i, m_i, s_i\}, \mathcal{A}_{\mathrm{drug}})
\end{equation}
where $\mathcal{A}_{\mathrm{drug}}$ parameterizes the installed aperture. The therapeutic effect is
\begin{equation}
  \boxed{\Delta \mathcal{S} = \mathcal{S}_{\mathrm{post}} - \mathcal{S}_{\mathrm{pre}}}
\end{equation}
as the primary pharmacological observable.
\end{theorem}

For an antidepressant transitioning a patient from turbulent ($\Rorder = 0.31$) to coherent ($\Rorder = 0.95$):
\begin{equation}
  \Delta \mathcal{S} = 10.26 - 0.965 = 9.30.
\end{equation}

This is a tenfold increase in the structural factor---a massive geometric reorganization of the neural circuit's partition structure.

\subsection{Three Mechanisms of Drug Action}

Drug action operates through three mechanisms, each modifying $\mathcal{S}$ differently.

\begin{definition}[Mechanism I: Aperture Installation]
\label{def:mech_aperture}
The drug installs a categorical aperture restricting the circuit's accessible phase space:
\begin{equation}
  \Omega_{\mathrm{neural}} \xrightarrow{\mathcal{A}_{\mathrm{drug}}} \Omega'_{\mathrm{neural}}, \quad \mu(\Omega') < \mu(\Omega).
\end{equation}
This directly modifies $\mathcal{S}$ through the aperture-dominated structural factor $\mathcal{S}_{\mathrm{aperture}} = n^2/n_{\max}^2$.
\end{definition}

\begin{definition}[Mechanism II: Coupling Modulation]
\label{def:mech_coupling}
The drug modulates the Kuramoto coupling strength:
\begin{equation}
  K(t) = K_{\mathrm{baseline}} + \Delta K_{\mathrm{drug}}(t),
\end{equation}
shifting the system toward or away from $K_c$. This modifies $\mathcal{S}$ through the order parameter: increased $K$ increases $\Rorder$, increasing $\mathcal{S}_{\mathrm{coherent}}$.
\end{definition}

\begin{definition}[Mechanism III: Regime Transition]
\label{def:mech_regime}
When coupling crosses $K_c$, the system undergoes a phase transition between regimes, discontinuously changing the functional form of $\mathcal{S}$:
\begin{equation}
  \mathcal{S}_{\mathrm{turbulent}} \xrightarrow{K > K_c} \mathcal{S}_{\mathrm{coherent}}.
\end{equation}
This is the therapeutic event.
\end{definition}

\subsection{Regime Transition Pathways by Drug Class}

Different drug classes induce characteristic pathways through the regime diagram.

\begin{theorem}[Drug Class Transition Pathways]
\label{thm:pathways}
The regime transition pathway depends on aperture type:
\begin{enumerate}[nosep]
  \item \textbf{SSRI (monopole):} Turbulent $\to$ Aperture-Dominated $\to$ Phase-Locked $\to$ Coherent. The single-target aperture first restricts accessible state space, then enables phase-locking as serotonin levels rise.
  \item \textbf{SNRI (dipole):} Turbulent $\to$ Hierarchical Cascade $\to$ Coherent. The dual-channel aperture (SERT + NET) creates hierarchical information flow bypassing the aperture-dominated stage.
  \item \textbf{TCA (quadrupole):} Turbulent $\to$ Phase-Locked $\to$ Coherent. The broad multi-target aperture directly enhances coupling across multiple receptor systems, driving rapid phase-locking.
  \item \textbf{Ketamine (NMDA antagonist):} Turbulent $\to$ transient Coherent. Rapid glutamatergic coupling enhancement crosses $K_c$ within hours rather than weeks~\citep{berman2000antidepressant,zarate2006randomized}.
  \item \textbf{Psychedelics:} Coherent $\to$ Turbulent $\to$ reorganized Coherent. Controlled regime destruction followed by reconstruction with modified partition geometry~\citep{carhart2016neural,petri2014homological}.
  \item \textbf{Anesthetics:} Any $\to$ Collapse ($\Rorder \to 0$, $\sigma^2 \to \sigma^2_{\max}$). Complete decoupling eliminates phase synchronization~\citep{alkire2008consciousness,mashour2005cognitive}.
\end{enumerate}
\end{theorem}

\subsection{Cross-Modal Therapeutic Equivalence}

\begin{theorem}[Cross-Modal Equivalence]
\label{thm:crossmodal}
Drugs that induce the same regime transition produce equivalent clinical outcomes regardless of binding profile:
\begin{equation}
  \Gstate_{\mathrm{SSRI}} = \Gstate_{\mathrm{TCA}} = \Gstate_{\mathrm{atypical}} = \Gstate_{\mathrm{coherent}}
\end{equation}
if all three reach the coherent regime from the turbulent regime.
\end{theorem}

\begin{proof}
The equation of state depends on $\mathcal{S}$, which depends on regime and partition geometry, not on the molecular identity of the aperture. Different apertures (monopole, dipole, quadrupole) modify different coordinates of partition space but converge on the same regime. The clinical outcome is determined by the final structural factor $\mathcal{S}_{\mathrm{coherent}}$, not by the transition pathway.
\end{proof}

This explains the convergent $\sim\!60\%$ response rate across antidepressant classes~\citep{cipriani2018comparative}: all achieve the turbulent-to-coherent regime transition, albeit through different pathways. The $\sim\!40\%$ non-response reflects patients in deep turbulent basins where no single aperture type provides sufficient $\Delta K$ to cross $K_c$.

\subsection{Onset Delay}

\begin{theorem}[Antidepressant Onset Delay]
\label{thm:onset}
The time to therapeutic response is
\begin{equation}
  \tau_{\mathrm{onset}} = \frac{1}{\gamma_{\mathrm{relax}}} \ln\left(\frac{K_c - K_{\mathrm{baseline}}}{K_c - K_{\mathrm{baseline}} - \Delta K_{\mathrm{drug}}}\right)
  \label{eq:onset}
\end{equation}
where $\gamma_{\mathrm{relax}}$ is the neural circuit's relaxation rate and $\Delta K_{\mathrm{drug}}$ is the drug-induced coupling enhancement.
\end{theorem}

\begin{proof}
The coupling strength evolves as $K(t) = K_{\mathrm{baseline}} + \Delta K_{\mathrm{drug}}(1 - e^{-\gamma_{\mathrm{relax}} t})$. Therapeutic response occurs when $K(t) \geq K_c$. Solving for $t$ gives Eq.~\eqref{eq:onset}.
\end{proof}

For typical parameters ($\gamma_{\mathrm{relax}} \sim 0.05$ day$^{-1}$, $K_c/K_{\mathrm{baseline}} \sim 2$, $\Delta K_{\mathrm{drug}}/K_{\mathrm{baseline}} \sim 1.5$), $\tau_{\mathrm{onset}} \approx 14$--$28$ days, consistent with the observed 2--4 week antidepressant latency~\citep{stahl2021essential}.

Ketamine's rapid onset ($\sim$hours) is explained by a much larger $\Delta K_{\mathrm{drug}}$ through glutamatergic coupling, crossing $K_c$ before exponential saturation~\citep{zarate2006randomized}.

% ============================================================================
\section{Frequency Matching and Timescale Hierarchy}
\label{sec:frequency}
% ============================================================================

\subsection{The Oscillatory Substrate}

The neural circuit's oscillatory substrate spans thirteen orders of magnitude in frequency.

\begin{definition}[Timescale Hierarchy]
\label{def:timescale}
The system exhibits a five-level timescale hierarchy:
\begin{equation}
  \tau_\rho \sim 10^{-15}\;\mathrm{s} \ll \tau_{\Hplus} \sim 10^{-14}\;\mathrm{s} \ll \tau_{\mathrm{config}} \sim 10^{-1}\;\mathrm{s} \ll \tau_{\mathrm{state}} \sim 10^{0}\;\mathrm{s} \ll \tau_{\mathrm{drug}} \sim 10^{3}\;\mathrm{s}
\end{equation}
spanning 18 orders of magnitude from charge redistribution to pharmacological intervention.
\end{definition}

The $\Hplus$ electromagnetic field oscillates at $\omega_{\Hplus} = 4.06 \times 10^{13}$ Hz, derived from proton transfer timescales $\tau_{\Hplus} \approx 25$ fs measured in enzyme active sites and mitochondrial membranes~\citep{marcus1993electron,klinman2013importance}. Molecular oxygen configurations transition at $\sim\!10$ Hz, setting the clock for intracellular categorical state changes~\citep{herzberg1950spectra}. Neural oscillations span the $\delta$-band (1--4 Hz) through $\gamma$-band (30--50 Hz)~\citep{buzsaki2004neuronal,fries2005mechanism}.

\begin{theorem}[Adiabatic Decoupling]
\label{thm:adiabatic}
Adjacent timescale levels are separated by at least $10^1$ in timescale, enabling adiabatic decoupling: fast variables provide slowly varying parameters for slow dynamics, while slow variables provide boundary conditions for fast dynamics~\citep{born1927quantentheorie}.
\end{theorem}

\subsection{Drug--Circuit Frequency Coupling}

\begin{definition}[Frequency Matching Condition]
\label{def:freq_match}
A drug couples to the neural circuit when its characteristic frequency $\omega_{\mathrm{drug}}$ matches an oscillatory hole frequency $\omega_{\mathrm{hole}}$ within tolerance:
\begin{equation}
  |\omega_{\mathrm{drug}} - \omega_{\mathrm{hole}}| < \Delta\omega_{\mathrm{thresh}} \sim 10^{11} \;\mathrm{Hz}.
\end{equation}
\end{definition}

An oscillatory hole is a functional absence in the phase-lock network where local coherence falls below threshold:
\begin{equation}
  \mathrm{Hole}_i = \{j \in \mathcal{V} : |\phi_j(t) - \Theta(t)| > \theta_{\mathrm{coh}}\},
\end{equation}
where $\theta_{\mathrm{coh}} \approx \pi/4$ defines the coherence boundary~\citep{buzsaki2006rhythms}.

\begin{theorem}[Gear Network Amplification]
\label{thm:gear}
Upon frequency matching, allosteric coupling activates therapeutic pathways through gear network transformation:
\begin{equation}
  \omega_{\mathrm{therapeutic}} = G_{\mathrm{pathway}} \times \omega_{\mathrm{drug}},
\end{equation}
where $G_{\mathrm{pathway}}$ is the gear ratio of the allosteric network connecting drug binding site to effector output. Typical biological gear ratios range from $G \sim 10^2$ to $10^4$~\citep{hartl2011molecular}.
\end{theorem}

\subsection{Categorical Distance and Catalytic Efficiency}

\begin{definition}[Categorical Distance]
\label{def:dcat}
The categorical distance between two states $\mathbf{q}_1$ and $\mathbf{q}_2$ in partition space is
\begin{equation}
  \dcat(\mathbf{q}_1, \mathbf{q}_2) = \sqrt{(n_1 - n_2)^2 + (\ell_1 - \ell_2)^2 + (m_1 - m_2)^2 + (s_1 - s_2)^2}.
\end{equation}
\end{definition}

\begin{theorem}[Diffusion-Limited Catalysis]
\label{thm:diffusion}
When the categorical distance between substrate and product is $\dcat = 1$ (adjacent categories), the catalytic rate approaches the diffusion limit:
\begin{equation}
  \dcat = 1 \implies k_{\mathrm{cat}}/K_M \geq 10^8 \;\mathrm{M}^{-1}\mathrm{s}^{-1}.
\end{equation}
Enzymes with $\dcat > 1$ are sub-diffusion~\citep{wolfenden2001depth,bar2011diffusion}.
\end{theorem}

This applies directly to drug metabolism: the rate at which a drug is processed depends on the categorical distance between its bound and unbound states in the enzyme's partition space.

% ============================================================================
\section{Pharmacological Trajectories in S-Entropy Space}
\label{sec:trajectories}
% ============================================================================

\subsection{The Trajectory-Terminus-Memory Triple}

A complete specification of the neural circuit's state requires not only the current configuration but also the path by which it arrived and the accumulated environmental changes along that path.

\begin{definition}[Circuit State Triple]
\label{def:triple}
The complete state of the circuit is
\begin{equation}
  \boxed{\Mstate = (\gamma, \Gstate_f, M)}
\end{equation}
where $\gamma: [0, T] \to \Sspace$ is the trajectory through S-entropy space, $\Gstate_f = \gamma(T)$ is the terminus (current configuration), and $M = \int_0^T (d\Hfield/dt)\,dt$ is the accumulated $\Hplus$ field change (memory).
\end{definition}

\begin{theorem}[Trajectory Non-Identity]
\label{thm:nonidentity}
If two trajectories $\gamma_1 \neq \gamma_2$ reach the same terminus $\Gstate_{f,1} = \Gstate_{f,2}$, the resulting circuit states are nonetheless distinct: $\Mstate_1 \neq \Mstate_2$.
\end{theorem}

This has immediate pharmacological significance: two drugs that produce the same receptor state (same terminus) via different binding pathways (different trajectories) produce different circuit states---different memories, different subsequent dynamics, different side effect profiles.

\subsection{Drug Action as Trajectory Modification}

\begin{theorem}[Pharmacological Trajectory Modulation]
\label{thm:drug_traj}
A drug modulates the circuit state by altering the trajectory while potentially preserving the terminus:
\begin{equation}
  \boxed{\Phi_{\mathrm{drug}}: (\gamma, \Gstate_f, M) \to (\gamma', \Gstate_f, M')}
\end{equation}
where $\gamma' \neq \gamma$ and $M' \neq M$, but $\Gstate_f$ may be unchanged.
\end{theorem}

This explains why anxiolytics reduce anxiety without eliminating threat assessment: they change the trajectory (how the thought of threat was reached) and the memory (accumulated emotional context) while leaving the configuration (threat detection capability) intact~\citep{stahl2021essential}.

\subsection{Variance Minimization as Therapeutic Mechanism}

\begin{theorem}[Free Energy--Variance Identity]
\label{thm:variance}
The Helmholtz free energy of a partition-based neural circuit is proportional to the phase variance:
\begin{equation}
  \boxed{F = \kB T \cdot \sigma^2(\phi)}
\end{equation}
The minimum achievable variance is
\begin{equation}
  \sigma^2_{\min} = \frac{\kB T}{K_{\mathrm{coupling}}},
\end{equation}
and the power consumed to maintain variance below thermal equilibrium is
\begin{equation}
  P = \gamma_{\mathrm{relax}} \cdot \kB T \cdot (\sigma^2_{\mathrm{thermal}} - \sigma^2).
\end{equation}
\end{theorem}

Drug action, thermodynamically, is variance reduction. A drug that increases $K_{\mathrm{coupling}}$ lowers $\sigma^2_{\min}$, enabling the circuit to achieve tighter phase synchronization---moving from turbulent to coherent operation.

\subsection{Poincar\'e Computing for Drug Design}

\begin{principle}[Poincar\'e Computing]
\label{princ:poincare}
A computation is specified by declaring the completion condition (terminus $\Gstate_f$) and constraint structure. The runtime determines the unique trajectory $\gamma^*$ that reaches $\Gstate_f$ while satisfying all constraints:
\begin{equation}
  \exists!\,[\gamma] : \gamma(T) = \Gstate_f \;\wedge\; \mathrm{Constraints}(\gamma) = \mathrm{true}.
\end{equation}
\end{principle}

Applied to drug design: specify the target regime (e.g., coherent with $\Rorder > 0.8$) and the partition structure determines the required aperture type, coupling modulation, and transition pathway. This is backward determination: specify WHERE the therapy must end, and the constraint structure determines HOW to get there.

% ============================================================================
\section{The Pharmacological Neural Partition Language}
\label{sec:pnpl}
% ============================================================================

We extend the Neural Partition Language (NPL) with pharmacological primitives grounded in the equations of state. The resulting pharmacological NPL (pNPL) provides a formal language for expressing drug action as operations on the neural circuit's partition structure.

\subsection{Type System}

\begin{definition}[pNPL Core Types]
\label{def:pnpl_types}
The pNPL type system extends NPL with pharmacological types:

\textbf{Circuit Types} (from partition structure):
\begin{itemize}[nosep]
  \item \texttt{PartitionCoord}: $(n, \ell, m, s)$ with capacity $C(n) = 2n^2$
  \item \texttt{SCoord}: $(\Sk, \St, \Se) \in [0,1]^3$
  \item \texttt{Regime}: $\{\text{coherent}, \text{turbulent}, \text{cascade}, \text{aperture}, \text{sync}\}$
  \item \texttt{StructuralFactor}: $\mathcal{S} \in \mathbb{R}^+$
\end{itemize}

\textbf{Pharmacological Types} (drug-specific):
\begin{itemize}[nosep]
  \item \texttt{Aperture}: $(\mathcal{A}, \text{type} \in \{\text{monopole}, \text{dipole}, \text{quadrupole}\}, \alpha)$
  \item \texttt{CouplingMod}: $\Delta K \in \mathbb{R}$
  \item \texttt{DrugTrajectory}: $\Phi_{\mathrm{drug}}: \Mstate \to \Mstate$
  \item \texttt{TransitionPath}: ordered sequence of \texttt{Regime}
\end{itemize}

\textbf{Observable Types} (measurable):
\begin{itemize}[nosep]
  \item \texttt{OrderParam}: $\Rorder \in [0, 1]$
  \item \texttt{PhaseVariance}: $\sigma^2 \in \mathbb{R}^+$
  \item \texttt{PLV}: phase-locking value $\in [0, 1]$
  \item \texttt{OnsetTime}: $\tau_{\mathrm{onset}} \in \mathbb{R}^+$
\end{itemize}
\end{definition}

\subsection{Core Operators}

\begin{definition}[pNPL Operators]
\label{def:pnpl_ops}

\textbf{Aperture Operations}:
\begin{itemize}[nosep]
  \item \texttt{INSTALL\_APERTURE}$(\Omega, \mathcal{A}_{\mathrm{drug}}) \to \Omega'$: Restrict phase space
  \item \texttt{CLASSIFY\_APERTURE}$(K_i\text{-profile}) \to \{\text{monopole}, \text{dipole}, \text{quadrupole}\}$
  \item \texttt{CATALYTIC\_FACTOR}$(\Omega, \Omega') \to \alpha$: Compute restriction ratio
\end{itemize}

\textbf{Regime Operations}:
\begin{itemize}[nosep]
  \item \texttt{REGIME\_CLASSIFY}$(\Rorder, \sigma^2) \to$ \texttt{Regime}: Determine current regime
  \item \texttt{STRUCTURAL\_FACTOR}$(\text{Regime}, \text{params}) \to \mathcal{S}$: Compute equation of state
  \item \texttt{DELTA\_S}$(\mathcal{S}_{\mathrm{pre}}, \mathcal{S}_{\mathrm{post}}) \to \Delta\mathcal{S}$: Therapeutic effect magnitude
\end{itemize}

\textbf{Transition Operations}:
\begin{itemize}[nosep]
  \item \texttt{TRANSITION\_PATH}$(\text{Regime}_{\mathrm{src}}, \text{Regime}_{\mathrm{tgt}}, \text{ApertureType}) \to$ \texttt{TransitionPath}
  \item \texttt{ONSET\_TIME}$(K_{\mathrm{base}}, \Delta K, K_c, \gamma_{\mathrm{relax}}) \to \tau_{\mathrm{onset}}$
  \item \texttt{CROSS\_MODAL\_TEST}$(\text{Drug}_1, \text{Drug}_2, \Gstate_{\mathrm{tgt}}) \to \text{bool}$
\end{itemize}

\textbf{Coupling Operations}:
\begin{itemize}[nosep]
  \item \texttt{MODULATE\_K}$(K_{\mathrm{base}}, \Delta K_{\mathrm{drug}}) \to K$: Coupling modulation
  \item \texttt{CRITICAL\_K}$(\sigma(\omega), \langle k \rangle) \to K_c$: Compute transition threshold
  \item \texttt{PHASE\_LOCK}$(\{\phi_i\}, K) \to \Rorder$: Kuramoto synchronization
\end{itemize}

\textbf{Frequency Operations}:
\begin{itemize}[nosep]
  \item \texttt{FREQ\_MATCH}$(\omega_{\mathrm{drug}}, \omega_{\mathrm{hole}}, \Delta\omega) \to \text{bool}$
  \item \texttt{GEAR\_ACTIVATE}$(\omega_{\mathrm{drug}}, G_{\mathrm{pathway}}) \to \omega_{\mathrm{therapeutic}}$
\end{itemize}

\textbf{Trajectory Operations}:
\begin{itemize}[nosep]
  \item \texttt{DRUG\_MODIFY}$(\Mstate, \Phi_{\mathrm{drug}}) \to \Mstate'$: Apply trajectory modification
  \item \texttt{VARIANCE\_MINIMIZE}$(\sigma^2, K, T) \to \sigma^2_{\min}$: Compute therapeutic floor
  \item \texttt{COMPLETE}$(\Gstate_f) \to \gamma^*$: Poincar\'e backward determination
\end{itemize}

\textbf{Measurement Operations}:
\begin{itemize}[nosep]
  \item \texttt{PLV\_MEASURE}$(\text{EEG channels}) \to \Rorder$: Experimental proxy
  \item \texttt{REGIME\_MAP}$(\Rorder(t)) \to \text{Regime}(t)$: Temporal regime tracking
  \item \texttt{CATEGORICAL\_THERM}$(\Delta\Se) \to T$: Zero-backaction thermometry
\end{itemize}
\end{definition}

\subsection{Program Structure: Declarative Therapeutic Completion}

pNPL programs are declarative: specify the therapeutic completion condition, and the partition structure determines the trajectory.

\begin{algorithm}
\caption{Antidepressant Treatment in pNPL}
\label{alg:antidepressant}
\begin{algorithmic}[1]
\State \textbf{THERAPEUTIC\_PROCESS} TreatDepression
\State
\State \Comment{Assess current regime}
\State $\Rorder_0 \gets \texttt{PLV\_MEASURE}(\text{patient EEG})$
\State $\text{regime}_0 \gets \texttt{REGIME\_CLASSIFY}(\Rorder_0, 1 - \Rorder_0)$
\State
\State \Comment{Declare completion: coherent regime}
\State $\Gstate_{\mathrm{target}} \gets \texttt{COMPLETE}(\text{Regime} = \text{coherent}, \Rorder > 0.8)$
\State
\State \Comment{Select aperture type from drug class}
\State $\mathcal{A} \gets \texttt{CLASSIFY\_APERTURE}(K_i\text{-profile})$
\State
\State \Comment{Compute transition pathway}
\State $\text{path} \gets \texttt{TRANSITION\_PATH}(\text{regime}_0, \text{coherent}, \mathcal{A})$
\State
\State \Comment{Predict onset time}
\State $K_c \gets \texttt{CRITICAL\_K}(\sigma(\omega), \langle k \rangle)$
\State $\tau \gets \texttt{ONSET\_TIME}(K_0, \Delta K_{\mathrm{drug}}, K_c, \gamma)$
\State
\State \Comment{Monitor regime transition}
\State \textbf{while} $\texttt{REGIME\_CLASSIFY}(\Rorder(t)) \neq \text{coherent}$ \textbf{do}
\State \quad $\Rorder(t) \gets \texttt{PLV\_MEASURE}(\text{patient EEG})$
\State \quad $\Delta\mathcal{S}(t) \gets \texttt{DELTA\_S}(\mathcal{S}_0, \texttt{STRUCTURAL\_FACTOR}(\Rorder(t)))$
\State \textbf{end while}
\State
\State \textbf{return} $\Mstate_{\mathrm{final}} = (\gamma^*, \Gstate_{\mathrm{target}}, M)$
\end{algorithmic}
\end{algorithm}

\begin{algorithm}
\caption{Drug Design via Poincar\'e Computing in pNPL}
\label{alg:drug_design}
\begin{algorithmic}[1]
\State \textbf{DESIGN\_PROCESS} TargetedDrug
\State
\State \Comment{Specify therapeutic target}
\State $\Gstate_f \gets (\Rorder > 0.8, \sigma^2 < 0.2, \text{Regime} = \text{coherent})$
\State
\State \Comment{Characterize patient's current state}
\State $\Mstate_0 \gets (\gamma_0, \Gstate_0, M_0)$ from clinical assessment
\State
\State \Comment{Compute required aperture properties}
\State $\Delta K_{\mathrm{required}} \gets K_c - K_0$
\State $\alpha_{\mathrm{required}} \gets \texttt{CATALYTIC\_FACTOR}(\Omega_0, \Omega_{\mathrm{target}})$
\State
\State \Comment{Determine aperture type}
\State \textbf{if} $\Delta K_{\mathrm{required}} < K_c / 3$ \textbf{then}
\State \quad type $\gets$ monopole \Comment{mild: single target suffices}
\State \textbf{else if} $\Delta K_{\mathrm{required}} < 2K_c / 3$ \textbf{then}
\State \quad type $\gets$ dipole \Comment{moderate: dual target}
\State \textbf{else}
\State \quad type $\gets$ quadrupole \Comment{severe: broad aperture}
\State \textbf{end if}
\State
\State \Comment{Backward determination of drug properties}
\State $\omega_{\mathrm{drug}} \gets \texttt{FREQ\_MATCH}^{-1}(\omega_{\mathrm{hole}}, \Delta\omega)$
\State $K_i\text{-profile} \gets \texttt{CLASSIFY\_APERTURE}^{-1}(\text{type})$
\State
\State \textbf{return} (type, $K_i$-profile, $\omega_{\mathrm{drug}}$, $\tau_{\mathrm{onset}}$)
\end{algorithmic}
\end{algorithm}

% ============================================================================
\section{Temporal Dynamics of Drug Action}
\label{sec:time}
% ============================================================================

\subsection{Circuit Completion and Subjective Time}

Subjective time is the duration of geometric tracing during circuit completion events~\citep{james1890principles,poppel1997temporal}.

\begin{theorem}[Internal Time]
\label{thm:time}
The subjective duration experienced during drug action equals the sum of circuit completion times for active oscillatory holes:
\begin{equation}
  T_{\mathrm{internal}} = \sum_{i \in \mathrm{active\;holes}} \tau_{\mathrm{circuit}}^{(i)}.
\end{equation}
\end{theorem}

\begin{corollary}[Specious Present]
The experiential ``now'' has duration
\begin{equation}
  \Delta t_{\mathrm{now}} \sim 100\text{--}1000\;\mathrm{ms},
\end{equation}
equal to the average circuit completion time for coherent oscillatory hole ensembles~\citep{poppel1997temporal,buhusi2005time}.
\end{corollary}

\subsection{Drug Effects on Time Perception}

\begin{theorem}[Temporal Elasticity Under Pharmacological Perturbation]
\label{thm:temporal_elasticity}
Drug-induced changes to circuit completion efficiency produce systematic time perception distortions:
\begin{equation}
  \frac{T_{\mathrm{subjective}}}{T_{\mathrm{objective}}} = \frac{\tau_{\mathrm{circuit, drug}}}{\tau_{\mathrm{circuit, baseline}}}.
\end{equation}
\end{theorem}

Specific predictions:
\begin{itemize}[nosep]
  \item \textbf{Stimulants} (increased electron transport efficiency): $\tau_{\mathrm{circuit}}$ decreases $\to$ time overestimation ($+15\% \pm 3\%$ predicted)
  \item \textbf{Depressants/hypoxia} (decreased transport): $\tau_{\mathrm{circuit}}$ increases $\to$ time underestimation ($-23\% \pm 4\%$ predicted)
  \item \textbf{Psychedelics}: Variable $\tau_{\mathrm{circuit}}$ due to regime disruption $\to$ profound time distortion~\citep{carhart2016neural}
  \item \textbf{Anesthetics}: $\tau_{\mathrm{circuit}} \to \infty$ as $\Rorder \to 0$ $\to$ subjective time cessation~\citep{alkire2008consciousness}
\end{itemize}

\subsection{Consciousness Window Modulation}

\begin{definition}[Consciousness as Decay Curve Intersection]
\label{def:consciousness}
Consciousness $\Ccat$ at time $t$ is the intersection of perception and thought decay curves:
\begin{equation}
  \Ccat(t) = \Pdecay(t) \cap \Tdecay(t),
\end{equation}
where $\Pdecay$ decays with time constant $\tau_P \approx 50$ ms and $\Tdecay$ with $\tau_T \approx 100$ ms. The consciousness window has characteristic duration:
\begin{equation}
  \Delta t_C = \frac{\tau_P \tau_T}{\tau_P + \tau_T} \ln\left(\frac{\Pdecay(0) \Tdecay(0)}{\theta^2}\right) \sim 100\text{--}500\;\mathrm{ms}.
\end{equation}
\end{definition}

\begin{theorem}[Pharmacological Window Modulation]
\label{thm:window}
Drug classes modify $\tau_P$ and $\tau_T$ differentially:
\begin{center}
\begin{tabular}{lccc}
\toprule
Drug Class & $\tau_P$ & $\tau_T$ & Window Effect \\
\midrule
Anxiolytics & $\uparrow$ & $\uparrow$ & Wider, calmer \\
Stimulants & $\downarrow$ & $\downarrow$ & Narrower, intense \\
Psychedelics & Variable & $\uparrow\uparrow$ & Distorted geometry \\
Anesthetics & $\to 0$ & $\to 0$ & No intersection \\
\bottomrule
\end{tabular}
\end{center}
\end{theorem}

% ============================================================================
\section{Experimental Predictions and Validation Pathways}
\label{sec:predictions}
% ============================================================================

The framework generates quantitative predictions testable with existing clinical and laboratory instrumentation.

\subsection{Molecular Scale Predictions}

\begin{enumerate}[nosep]
  \item \textbf{Proton flux frequency}: $\omega_{\Hplus} = 4.06 \times 10^{13}$ Hz from partition structure, testable via ultrafast infrared spectroscopy of mitochondrial membranes~\citep{mitchell1961coupling,wikstrom2015proton}.
  \item \textbf{Categorical distance--efficiency correlation}: $\dcat$ vs $\log_{10}(k_{\mathrm{cat}}/K_M)$ should yield $r < -0.8$ across enzyme classes.
  \item \textbf{Adiabatic separation}: $\omega_{\Hplus}/\omega_{\mathrm{config}} \sim 10^{12}$, confirming timescale hierarchy.
\end{enumerate}

\subsection{Neural Scale Predictions}

\begin{enumerate}[nosep]
  \item \textbf{PLV--Kuramoto correspondence}: PLV measured from clinical EEG should correlate with the Kuramoto order parameter with $r > 0.95$~\citep{lachaux1999measuring}.
  \item \textbf{Depression as turbulent regime}: Depressed patients should show $\Rorder < 0.3$ in resting-state $\alpha$-band EEG; healthy controls $\Rorder > 0.8$~\citep{fingelkurts2006impaired}.
  \item \textbf{Sleep stage regime mapping}: Waking = coherent, N3 = phase-locked, REM = turbulent, with $\sigma^2_{\mathrm{REM}} > \sigma^2_{\mathrm{N3}}$~\citep{steriade2003neuronal,diekelmann2010memory}.
  \item \textbf{Poincar\'e deviation}: Zero exact returns ($\delta < 10^{-10}$) in bounded trajectory analysis over $> 100$ time windows.
\end{enumerate}

\subsection{Pharmacological Scale Predictions}

\begin{enumerate}[nosep]
  \item \textbf{Cross-modal equivalence}: Drugs achieving the same regime transition should produce equivalent clinical response rates ($\sim\!60\%$) regardless of binding profile distance~\citep{cipriani2018comparative}.
  \item \textbf{Regime transition monitoring}: Longitudinal EEG during antidepressant treatment should show progressive $\Rorder$ increase from $\sim\!0.3$ to $\sim\!0.9$ over 2--4 weeks, with regime transitions visible as discontinuities in $\mathcal{S}(t)$.
  \item \textbf{Onset delay prediction}: Eq.~\eqref{eq:onset} should predict $\tau_{\mathrm{onset}}$ to within $\pm 5$ days for SSRIs, SNRIs, and TCAs when $K_{\mathrm{baseline}}$, $\Delta K_{\mathrm{drug}}$, and $K_c$ are estimated from pre-treatment EEG.
  \item \textbf{Treatment resistance prediction}: Patients with $K_{\mathrm{baseline}} \ll K_c$ (deep turbulent basin) should be identifiable pre-treatment as non-responders to monotherapy~\citep{rush2006acute}.
  \item \textbf{Dose--response as structural factor curve}: Dose--response curves should map to $\Delta\mathcal{S}(\text{dose})$ with sigmoidal shape reflecting the critical coupling transition.
  \item \textbf{Ketamine rapid onset}: Ketamine's onset time from Eq.~\eqref{eq:onset} should be $< 24$ hours due to large glutamatergic $\Delta K$~\citep{zarate2006randomized}.
\end{enumerate}

\subsection{Temporal Predictions}

\begin{enumerate}[nosep]
  \item \textbf{Circuit completion--duration correlation}: Measured $\tau_{\mathrm{circuit}}$ (from EEG oscillatory analysis) should correlate with subjective duration estimates with $r > 0.85$.
  \item \textbf{Time perception under pharmacological manipulation}: Electron transport inhibitors should produce systematic duration underestimation; enhancers should produce overestimation.
\end{enumerate}

\subsection{Proposed Experimental Protocol}

\begin{enumerate}
  \item \textbf{Longitudinal regime tracking study}: $N = 60$ patients with major depression, randomized to SSRI, SNRI, or TCA. Weekly 19-channel resting-state EEG with PLV analysis. Primary endpoint: $\Rorder(t)$ trajectory and regime transition timing vs.\ Eq.~\eqref{eq:onset} prediction.

  \item \textbf{Structural factor dose--response}: $N = 30$ healthy volunteers. Single ascending dose of a well-characterized SSRI with pre- and post-dose EEG at 1, 4, 8, 24 hours. Primary endpoint: $\Delta\mathcal{S}(\text{dose})$ curve shape.

  \item \textbf{Cross-modal equivalence test}: $N = 120$ patients, randomized to SSRI, SNRI, TCA, or atypical. EEG regime classification at weeks 0, 2, 4, 8. Primary endpoint: convergence of $\Rorder_{\mathrm{final}}$ across drug classes among responders.

  \item \textbf{Treatment resistance prediction}: $N = 100$ treatment-na\"ive patients. Baseline EEG for $K_{\mathrm{baseline}}$ estimation. Prediction of responder/non-responder status before treatment initiation. Validation against 8-week clinical outcomes.

  \item \textbf{Time perception pharmacology}: $N = 40$ healthy volunteers in crossover design (drug vs.\ placebo). Temporal bisection and duration reproduction tasks under SSRI, stimulant, and anxiolytic conditions. Primary endpoint: $T_{\mathrm{subjective}}/T_{\mathrm{objective}}$ ratio vs.\ $\tau_{\mathrm{circuit}}$ measured from concurrent EEG.
\end{enumerate}

% ============================================================================
\section{Clinical Implications}
\label{sec:clinical}
% ============================================================================

\subsection{Depression as Thermodynamic State}

The identification of depression with the turbulent operational regime ($\Rorder < 0.3$, $\mathcal{S}_{\mathrm{turbulent}} \approx 0.97$) provides a fundamentally different conceptualization from the ``chemical imbalance'' hypothesis~\citep{owens1997role}. Depression is not a deficiency of a neurotransmitter; it is a thermodynamic state of the neural oscillator network characterized by high phase variance and low inter-regional coherence~\citep{uhlhaas2006neural}.

This reconceptualization explains several otherwise puzzling clinical observations:

\begin{enumerate}[nosep]
  \item \textbf{Multiple effective drug classes}: All are trajectory modifiers inducing the turbulent-to-coherent transition via different regime pathways~\citep{cipriani2018comparative}.
  \item \textbf{Delayed onset}: The 2--4 week delay is the time to cross $K_c$ (Theorem~\ref{thm:onset}).
  \item \textbf{Treatment resistance}: Deep turbulent basins where $\Delta K_{\mathrm{drug}} < K_c - K_{\mathrm{baseline}}$ for any single agent~\citep{rush2006acute}.
  \item \textbf{Placebo response}: The $\sim\!30\%$ placebo rate reflects spontaneous regime transitions driven by expectation-mediated coupling enhancement.
  \item \textbf{Relapse}: Return to turbulent regime when $K$ drops below $K_c$ after drug withdrawal.
\end{enumerate}

\subsection{Combination Therapy as Parallel Aperture Installation}

\begin{theorem}[Combination Therapy]
\label{thm:combination}
Two drugs with orthogonal aperture types provide additive coupling enhancement:
\begin{equation}
  \Delta K_{\mathrm{combined}} = \Delta K_1 + \Delta K_2 - \Delta K_1 \Delta K_2 / K_{\max}.
\end{equation}
Combination therapy is predicted to be effective when $\Delta K_1 < K_c - K_{\mathrm{baseline}}$ and $\Delta K_2 < K_c - K_{\mathrm{baseline}}$ but $\Delta K_{\mathrm{combined}} > K_c - K_{\mathrm{baseline}}$.
\end{theorem}

This provides a principled basis for augmentation strategies in treatment-resistant depression: select drugs with orthogonal aperture types (e.g., SSRI monopole + lithium coupling enhancer) to maximize $\Delta K_{\mathrm{combined}}$.

\subsection{Personalized Medicine via Regime Mapping}

Pre-treatment EEG provides three quantities that predict treatment response:
\begin{enumerate}[nosep]
  \item $\Rorder_{\mathrm{baseline}}$: Current phase coherence (how deep in the turbulent basin)
  \item $K_c$: Critical threshold (estimated from frequency variance $\sigma(\omega)$ and network connectivity $\langle k \rangle$)
  \item $\Delta K_{\mathrm{required}} = K_c - K_{\mathrm{baseline}}$: Required coupling enhancement
\end{enumerate}

Patients with small $\Delta K_{\mathrm{required}}$ (shallow turbulent basin) respond to any drug class; patients with large $\Delta K_{\mathrm{required}}$ (deep basin) require quadrupole apertures or combination therapy. This enables treatment selection before exposure.

\subsection{Extension to Other Psychiatric Conditions}

The regime framework extends beyond depression:
\begin{itemize}[nosep]
  \item \textbf{Anxiety}: Aperture-dominated regime with excessive phase space restriction~\citep{uhlhaas2006neural}
  \item \textbf{Schizophrenia}: Chimera state---coexistence of coherent and turbulent subpopulations~\citep{panaggio2015chimera,uhlhaas2010neural}
  \item \textbf{Mania}: Hyper-coherent regime ($\Rorder \to 1$) with excessive synchronization
  \item \textbf{Epilepsy}: Pathological phase-locking exceeding functional range
  \item \textbf{ADHD}: Unstable cascade regime with frequent transitions
\end{itemize}

Each condition maps to a specific region of the $(\Rorder, \sigma^2)$ phase diagram, and pharmacotherapy corresponds to navigation toward the coherent regime.

% ============================================================================
\section{Discussion}
\label{sec:discussion}
% ============================================================================

\subsection{Relation to Existing Pharmacological Theory}

The framework subsumes rather than contradicts existing pharmacology. Receptor binding ($K_i$) remains relevant as the mechanism by which an aperture is installed, but the therapeutic effect is determined by the resulting structural factor change $\Delta\mathcal{S}$, not by binding affinity alone. This resolves the paradox of high-affinity drugs with poor efficacy and promiscuous binders with broad therapeutic utility~\citep{roth2004multiplicity}: what matters is the aperture's effect on circuit geometry, not the strength of molecular adhesion.

The free energy principle~\citep{friston2010free} is encompassed as a special case: variance minimization (Theorem~\ref{thm:variance}) is the thermodynamic basis for free energy minimization, derived here from the triple equivalence rather than postulated as a variational principle.

\subsection{Why Apertures, Not Demons}

Prior frameworks for biological information processing invoke Maxwell demon mechanisms~\citep{sagawa2010generalized}---measurement, feedback, and erasure cycles. The categorical aperture framework demonstrates that biological selectivity operates through topological constraint at zero thermodynamic work, not through information acquisition and erasure. The Szil\'ard--Landauer--Bennett argument~\citep{szilard1929entropieverminderung,landauer1961irreversibility,bennett1982thermodynamics} is not violated; it is inapplicable, because its premise---that the sorting agent acquires information---does not hold for apertures.

This distinction is not semantic but thermodynamic. It means that the extraordinary selectivity of ion channels, enzyme active sites, and drug--receptor interactions imposes no Landauer cost~\citep{doyle1998structure,hille2001ion}. These structures are apertures in electromagnetic field configurations, not information-processing demons.

\subsection{The Triple Equivalence as Physical Law}

The triple equivalence $\Sosc = \Scat = \Spart = \kB M \ln n$ (Theorem~\ref{thm:triple}) is not an analogy but a mathematical identity: three descriptions of the same bounded phase space yield the same entropy because they enumerate the same microstates. The experimental predictions in Section~\ref{sec:predictions} test whether biological neural circuits are well described by the three axioms---and therefore whether the triple equivalence holds as a physical law governing drug action.

\subsection{Limitations}

Several limitations merit acknowledgment. First, the regime classification thresholds ($\Rorder > 0.8$ for coherent, $\Rorder < 0.3$ for turbulent) are derived from theoretical considerations and require empirical calibration with clinical populations. Second, the onset delay formula (Eq.~\ref{eq:onset}) assumes exponential coupling buildup, which may not capture the full complexity of receptor adaptation and downstream signaling cascades. Third, the aperture taxonomy (monopole/dipole/quadrupole) is based on the number of high-affinity targets, which is sensitive to the $K_i$ threshold chosen ($< 100$ nM). Fourth, the variance minimization identity (Theorem~\ref{thm:variance}) is derived for equilibrium conditions; neural circuits operate far from equilibrium, and corrections may be required.

These limitations define the experimental program: each can be resolved through the validation studies proposed in Section~\ref{sec:predictions}.

% ============================================================================
\section{Conclusion}
\label{sec:conclusion}
% ============================================================================

We have derived pharmacological equations of state for hybrid microfluidic circuits from three axioms---bounded phase space, the No Null State principle, and finite observational resolution. The principal results are:

\begin{enumerate}
  \item \textbf{Drug action is structural factor modification.} The universal equation of state $PV = N\kB T \cdot \mathcal{S}$ governs neural circuit dynamics. Drugs modify the temperature-independent structural factor $\mathcal{S}$, not the thermal energy scale. Therapeutic effect is $\Delta\mathcal{S} = \mathcal{S}_{\mathrm{post}} - \mathcal{S}_{\mathrm{pre}}$.

  \item \textbf{Categorical apertures provide zero-cost selectivity.} Drug molecules are exogenous categorical apertures---topological constraints on phase space operating at zero thermodynamic work ($W_{\mathrm{aperture}} = 0$). The Maxwell demon paradox is resolved: biological selectivity is topological, not informational.

  \item \textbf{Five regimes constitute the pharmacological phase diagram.} Coherent ($\Rorder > 0.8$), turbulent ($\Rorder < 0.3$), hierarchical cascade, aperture-dominated, and phase-locked regimes each possess distinct equations of state. Drug classes induce characteristic regime transition pathways.

  \item \textbf{Cross-modal equivalence explains convergent outcomes.} Chemically diverse drugs that achieve the same regime transition produce equivalent clinical outcomes, explaining the convergent $\sim\!60\%$ antidepressant response rate across drug classes.

  \item \textbf{Onset delay derives from critical coupling dynamics.} The 2--4 week antidepressant latency equals the time for sustained coupling enhancement to cross $K_c$. Ketamine's rapid onset reflects large glutamatergic $\Delta K$.

  \item \textbf{Variance minimization is the therapeutic mechanism.} $F = \kB T \cdot \sigma^2(\phi)$ identifies drug action with phase variance reduction toward $\sigma^2_{\min} = \kB T / K_{\mathrm{coupling}}$.

  \item \textbf{The pharmacological NPL provides computational primitives.} Aperture operations, regime transitions, frequency matching, and trajectory modification operators enable formal expression of drug action as constraint satisfaction in partition space.

  \item \textbf{Temporal dynamics derive from circuit completion.} Drug effects on time perception, consciousness window duration, and subjective experience follow from circuit completion mechanics modulated by pharmacological perturbation.
\end{enumerate}

The framework does not merely draw analogies between thermodynamics and pharmacology: it proves, from physically realized axioms, that the same entropy governs molecular transport, neural oscillation, and clinical drug response. Drug design becomes categorical state navigation in S-entropy space. Treatment selection becomes regime transition pathway optimization. Clinical monitoring becomes real-time structural factor tracking.

The experimental program defined in Section~\ref{sec:predictions} provides the blueprint for validating these predictions. If the neural circuit satisfies bounded phase space, the No Null State principle, and finite observational resolution---as all physical systems must---then the pharmacological equations of state derived here govern drug action as a consequence of thermodynamic law.

% ============================================================================
\section*{Acknowledgments}
% ============================================================================

The author acknowledges the KEGG database~\citep{kanehisa2000kegg}, ChEMBL database~\citep{gaulton2012chembl}, OpenNeuro platform~\citep{markiewicz2021openneuro}, and PhysioNet archive~\citep{goldberger2000physiobank} for providing open access to data informing the validation framework.

% ============================================================================
\bibliographystyle{plainnat}
\bibliography{references}

\end{document}
